\paragraph{Bilingual supervision changes which facts a glyph writer returns.}
On the same pictured Chinese novel passages, a writer trained with aligned English questions gave 121 exact Chinese answer images to 748 English requests. Giving the same English questions to the same passages during training but associating each with a different fact yielded 54 exact replies. Repeating the Chinese training question yielded seven, and the parent gave ten. For Chinese questions the corresponding counts were 232, 228, 213 and 93. On the first passage with four exact stopped replies, both question languages elicited the correct location (the kitchen) and the correct person (Daisy). Elsewhere, an English question about rent elicited a sea view, while the Chinese form elicited the recorded eighty dollars (Fig.~\ref{fig:bilingual-novel}).

Three of 374 novel passages and one of 227 regular-print XQuAD passages received all four exact answers from the aligned model. Both training controls completed zero passages on both tests. The gain of three novel passages had a paired 95\% bootstrap interval of 0--1.87 percentage points versus the exchanged control (two-sided paired sign test, $p=0.25$). The aligned model retained both regular-print CMRC replies on 16 of 197 passages and both familiar historical requests on 99 of 100 headwords. It met the predetermined selection rule for the bilingual research checkpoint. It completed no four-request passage on XQuAD in HanaMin and gave zero canonical replies to 300 questions requiring content ink absent from the input.

Wrong answers generally began with the wrong glyph. Of 627 missed English requests, 573 diverged from every accepted reference at the first cell and 567 returned a contiguous excerpt from the source. The corresponding counts for Chinese were 428 and 447 of 516. Every non-whitespace English input cell was already in the optical inventory. Among the 462 questions for which none of the annotated answers had appeared as a supervised target in our BiPaR fit, the writer gave 40 exact English-request and 102 exact Chinese-request replies. The evaluation passages were held out, although their novel titles occurred in the training partition. In a separate two-household example, exchanging dictionary counts changed a correctly written reply from three dictionaries to five. Changing the household named in the Chinese question on either page left that reply unchanged.

The aligned fit updated 28.60 million of 87.09 million network parameters for 2,496 steps on one RTX 4090 D. Optimization and checkpoint writing took 175.5 seconds with 4.50 GiB peak allocated GPU memory. The three fits had 538,340 supervised next actions each. Warm serial serving on the first 64 validation passages gave 42/128 exact Chinese-request and 19/128 English-request replies in 5.97 and 3.63 seconds. Local text Qwen3-8B Q8 gave 74/128 and 63/128 in 29.00 and 29.88 seconds on the same GPU. Qwen completed all four requests on 18 of these 64 passages and the image writer on zero. These are inference costs at different accuracies. The visual writer also inherits text pretraining whose cost is excluded from the measured adaptation.

At the first action that distinguished the two annotated answers, the aligned writer gave the intended glyph higher probability than its paired rival for 603/748 Chinese and 494/748 English requests. It ranked that glyph highest among all available actions for only 278 and 152 requests, respectively. A third action ranked highest for 440 Chinese and 536 English requests. Recomputing the selected action agreed with every saved autonomous decision for which the writer had followed the shared reference prefix (741 Chinese and 742 English). Ranking a fact ahead of the other annotated fact therefore overestimated the writer's accuracy when it had to choose freely from all glyphs.
